\chapter{The polar Neumann complex}
\label{mht:ch:polar-neumann}

A boundary condition can leave a state even when the unconstrained
Dolbeault complex has no first cohomology. On the sphere, every smooth
$(0,1)$-form has a primitive. If the primitive must have zero first
holomorphic derivative at a marked point, one scalar obstruction remains.
The meromorphic form with a double pole at that point measures precisely
this obstruction. We compute the resulting algebra and retain the trace
needed for its classical action.

\section{A first jet and its boundary integral}

Let $C=\mathbb P^1$ have its complex orientation, and let $o=0$ in the
affine coordinate $z$. Fix $\kappa\in\mathbb C^\times$ and put
\[
 \alpha=\kappa z^{-2}dz,\qquad j(f)=\partial_z f(o)
 \quad\text{for }f\in C^\infty(C).
\]
The form $\alpha$ extends across infinity and has a single double pole,
with zero residue. All smooth function and form spaces below carry their
usual Fr\'echet topologies.

\begin{definition}\label{mht:def:neumann-algebra}
The Neumann Dolbeault algebra is the cohomologically graded algebra
\[
 E^0=\ker j,\qquad E^1=\Omega^{0,1}(C),\qquad
 E^k=0\ (k\ne0,1),\qquad d_E=\bar\partial,
\]
with pointwise multiplication and wedge product.
\end{definition}

The condition is imposed on functions. It is preserved by multiplication,
since $j(fg)=f(o)j(g)+j(f)g(o)$. A product of two degree-one elements
vanishes. Thus $E$ is a unital commutative differential graded algebra.
Its holomorphic germs at $o$ form
$\mathbb C\oplus z^2\mathbb C\{z\}$, whose completion is
$\mathbb C\oplus z^2\mathbb C[[z]]$.

The polar form defines a trace by deleting round disks in the fixed
coordinate:
\begin{equation}\label{mht:eq:polar-trace}
 \tau_E(\eta)=\frac1{2\pi i}
   \lim_{r\downarrow0}\int_{C\setminus\{|z|<r\}}\alpha\wedge\eta
 \quad(\eta\in E^1),\qquad \tau_E(E^0)=0.
\end{equation}
The coordinate and the circular prescription are part of this definition.

\begin{lemma}[Polar Stokes identity]\label{mht:lem:polar-stokes}
The limit in \eqref{mht:eq:polar-trace} exists and defines a continuous
functional of degree $-1$. For every $f\in C^\infty(C)$,
\begin{equation}\label{mht:eq:polar-stokes}
 \tau_E(\bar\partial f)=\kappa j(f).
\end{equation}
In particular, $\tau_Ed_E=0$ on $E$, and
$\tau_E(xy)=(-1)^{|x||y|}\tau_E(yx)$ for homogeneous $x,y\in E$.
\end{lemma}

\begin{proof}
The vanishing of $H^{0,1}(\mathbb P^1)$ gives a continuous solution
operator
\begin{equation}\label{mht:eq:normalized-green}
 G:\Omega^{0,1}(C)\longrightarrow C^\infty(C),\qquad
 \bar\partial G\eta=\eta,\qquad (G\eta)(o)=0.
\end{equation}
One obtains it from a Dolbeault Green operator by subtracting the value
at $o$. The normalization makes it unique, because the global
holomorphic functions on $C$ are constant.

Away from $o$, one has $d(\alpha f)=-\alpha\wedge\bar\partial f$.
The inner boundary of the deleted disk is oriented clockwise, so
Stokes' theorem gives
\[
 \frac1{2\pi i}\int_{C\setminus\{|z|<r\}}\alpha\wedge\bar\partial f
 =\frac1{2\pi i}\oint_{|z|=r}^{\mathrm{counterclockwise}}
        \kappa z^{-2}f\,dz.
\]
In the expansion
$f=f(o)+zj(f)+\bar z\,\partial_{\bar z}f(o)+O(|z|^2)$,
only the holomorphic linear term contributes in the limit.
This proves \eqref{mht:eq:polar-stokes}.

Every $\eta$ is $\bar\partial G\eta$, so the same identity proves
existence and continuity of the
trace on all of $E^1$. It vanishes on $d_EE^0$ by the defining condition
$j=0$. Graded commutativity gives cyclicity.
\end{proof}

The trace is therefore the obstruction to finding a Neumann primitive.
To exhibit its class, choose a smooth function $f_0$ with
$f_0(o)=0$ and $j(f_0)=1$, and set $\eta_0=\bar\partial f_0$.
A useful choice, used in the explicit calculations below, is
\begin{equation}\label{mht:eq:sphere-representative}
 f_0=\frac{z}{1+z\bar z},\qquad
 \eta_0=-\frac{z^2}{(1+z\bar z)^2}\,d\bar z.
\end{equation}
In the coordinate $w=1/z$, the function is
$\bar w/(1+w\bar w)$, so it is smooth at infinity.
Although $\eta_0$ is exact in the unconstrained Dolbeault complex,
it is not exact in $E$, since $\tau_E(\eta_0)=\kappa$.

\section{The two-dimensional state algebra}

Let
\begin{equation}\label{mht:eq:residual-algebra}
 A_\kappa=\mathbb C\oplus\mathbb C\varepsilon,\qquad
 |\varepsilon|=1,\quad \varepsilon^2=0,\quad d_A=0,\qquad
 \tau_A(1)=0,\quad\tau_A(\varepsilon)=\kappa.
\end{equation}
Its inclusion into the smooth complex sends the odd generator to the
class just constructed:
\[
 i:A_\kappa\longrightarrow E,\qquad i(1)=1,\quad i(\varepsilon)=\eta_0.
\]

\begin{proposition}\label{mht:prop:polar-model}
The map $i$ is a continuous unital differential graded algebra
quasi-isomorphism, and $\tau_Ei=\tau_A$.
\end{proposition}

\begin{proof}
It preserves the differential and product because $\eta_0$ is closed
and $\eta_0^2=0$. Define
\[
 q(f)=f(o),\qquad q(\eta)=j(G\eta)\varepsilon,\qquad
 h_0(\eta)=G\eta-j(G\eta)f_0,\qquad h_0(E^0)=0.
\]
The first jet of $h_0\eta$ is zero.

For $f\in E^0$, the normalization
of $G$ gives $Gd_Ef=f-f(o)$, hence
$h_0d_Ef=f-iq(f)$ and $qd_Ef=0$. On degree one,
$d_Eh_0\eta=\eta-iq(\eta)$. Also $G\eta_0=f_0$, so
$qi=1$, $h_0i=qh_0=h_0^2=0$. These formulas give a continuous
contraction of complexes and prove the quasi-isomorphism.
The polar Stokes identity gives $\tau_Ei=\tau_A$.
\end{proof}

Point evaluation gives the simplest contraction, but an action also
uses the bilinear pairing $\langle x,y\rangle_E=\tau_E(xy)$.
For transfer of that pairing, the projection must be adjoint to $i$.
This requirement changes its degree-zero component.

\begin{theorem}[Cyclic contraction]\label{mht:thm:cyclic-contraction}
Define $p(f)=p_0(f)1$ and $p(\eta)=p_1(\eta)\varepsilon$ by
\begin{equation}\label{mht:eq:adjoint-projection}
 p_0(f)=\kappa^{-1}\tau_E(f\eta_0),\qquad
 p_1(\eta)=\kappa^{-1}\tau_E(\eta),
\end{equation}
and define a degree-$-1$ map by
\begin{equation}\label{mht:eq:cyclic-homotopy}
 h\eta=G\eta-p_1(\eta)f_0+\kappa^{-1}\tau_E(f_0\eta)1,
 \qquad h(E^0)=0.
\end{equation}
These continuous maps satisfy
\begin{equation}\label{mht:eq:contraction-identities}
 pi=1,\qquad d_Eh+hd_E=1-ip,\qquad h^2=hi=ph=0.
\end{equation}
The maps $p$ and $i$ are adjoint for $\tau_E(xy)$ and $\tau_A(xy)$.
The homotopy satisfies
\begin{equation}\label{mht:eq:cyclicity}
 \tau_E((h\eta)\xi)+\tau_E(\eta h\xi)=0
 \qquad(\eta,\xi\in E^1).
\end{equation}
\end{theorem}

\begin{proof}
By \eqref{mht:eq:polar-stokes}, $j(G\eta)=p_1(\eta)$, so $h\eta\in E^0$.
For $f\in E^0$, the same identity applied to $ff_0$ yields
\[
 \tau_E(f\eta_0)+\tau_E(f_0d_Ef)=\kappa f(o).
\]
Together with $Gd_Ef=f-f(o)$, this proves $hd_Ef=f-p_0(f)$.
Differentiating \eqref{mht:eq:cyclic-homotopy} gives
$d_Eh\eta=\eta-p_1(\eta)\eta_0$.

One has $p_0(1)=p_1(\eta_0)=1$ and
$\tau_E(f_0\eta_0)=\frac12\tau_Ed_E(f_0^2)=0$.
Consequently $pi=1$ and $hi=0$. Applying the polar Stokes identity to
$(G\eta)f_0$, whose value and first derivative at $o$ vanish, gives
\[
 \tau_E((G\eta)\eta_0)=-\tau_E(f_0\eta).
\]
This proves $ph=0$. Degree considerations give $h^2=0$.
The two formulas in \eqref{mht:eq:adjoint-projection} are exactly the
adjointness identities.

For cyclicity, set $F=G\eta$, $K=G\xi$,
$r=p_1(\eta)$, $s=p_1(\xi)$,
$u=\kappa^{-1}\tau_E(f_0\eta)$ and
$v=\kappa^{-1}\tau_E(f_0\xi)$. The left side of
\eqref{mht:eq:cyclicity} becomes
\[
 \tau_E(F\xi+\eta K)-\kappa rv-\kappa su+\kappa us+\kappa vr.
\]
The last four terms cancel. The first is
$\tau_Ed_E(FK)=\kappa j(FK)=0$, since $F(o)=K(o)=0$.
Continuity follows from that of $G$, $\tau_E$, and multiplication by
the fixed smooth function $f_0$.
\end{proof}

\begin{example}\label{mht:ex:projection-not-evaluation}
For the choice \eqref{mht:eq:sphere-representative}, put
$v=z\bar z/(1+z\bar z)$. Then $v\in E^0$ and $v(o)=0$, whereas
\[
 G(v\eta_0)=\frac{z}{1+z\bar z}
                -\frac{z}{2(1+z\bar z)^2},\qquad
 p_0(v)=\frac12.
\]
The displayed primitive is smooth in both charts, vanishes at $o$,
and differentiates to $v\eta_0$. Its first jet computes $p_0(v)$.
Thus point evaluation and the cyclic projection impose different
normalizations on the eliminated functions.
\end{example}

The contraction has a particularly simple effect on the internal
algebra. In the binary-tree formulas for transfer, leaves are labeled
by $i$, vertices by multiplication, internal edges by $h$, and the root
by $p$. Every tree with at least three leaves has a non-root vertex
joining two leaves. Their product lies in $i(A_\kappa)$ and is killed by
the following $h$. Hence all transferred products of arity at least
three vanish. The binary product is the one in
\eqref{mht:eq:residual-algebra}.

This statement concerns the internal coefficient algebra. It supplies
neither an inverse pairing on arbitrary distributional covectors nor
operations at coincident field insertions.

\section{Choice of representative and surface fields}

Changing $f_0$ changes the representative of the odd state, but leaves
its trace fixed. The comparison can be made over parameter simplices,
so the choices admit compatible homotopies.

\begin{proposition}\label{mht:prop:representative-homotopies}
An affine family $f$ of smooth functions satisfying $f(o)=0$ and
$j(f)=1$, parametrized by a simplex $P$, defines a trace-preserving
unital differential graded algebra map
\[
 A_\kappa\longrightarrow\Omega^\bullet_{\mathrm{poly}}(P)\otimes E,
 \qquad \varepsilon\longmapsto(d_P+\bar\partial)f.
\]
The maps commute with face and degeneracy restrictions.
\end{proposition}

\begin{proof}
Use the indicated tensor order and its Koszul signs. The coefficients
of $d_Pf$ have zero first jet, so they lie in $E^0$. The image of
$\varepsilon$ is closed and odd, so its square is zero.
The trace annihilates its $E^0$ component and sends
$\bar\partial f$ to $\kappa$. Restriction of the affine family commutes
with the formula.

On an interval, writing $u=f_1-f_0$, the image is
$\eta_0+t\bar\partial u+dt\,u$. The two differential terms
$dt\,\bar\partial u$ cancel.
\end{proof}

Let $U$ be an oriented smooth surface and $N\geq1$ an integer.
Tensor the comparison with $\Omega^\bullet(U)$ and
$\operatorname{Mat}_N(\mathbb C)$, using the completed projective tensor
product of the smooth spaces and the graded tensor product of algebras.
The homotopy on a simple tensor is
\[
 h_U(\omega\otimes\eta)=(-1)^{|\omega|}\omega\otimes h\eta.
\]
The identities \eqref{mht:eq:contraction-identities} then hold for the
total differential. These maps commute with restriction in $U$ and
preserve support there.

For compact support, take the inductive limit
over fixed compact subsets of $U$. The Green operator acts on the whole
compact curve, so it is not local in the curve variable.

The trace also determines the classical action. Let
\[
 L_U=\Omega^\bullet(U)\widehat\otimes E
                         \otimes\operatorname{Mat}_N(\mathbb C)
\]
with total form degree and the graded commutator bracket.

A physical connection is an element $\mathcal A\in L_U^1$.
Its surface component has Neumann function coefficients, while its
curve component is an arbitrary smooth $(0,1)$-form. Define
\begin{equation}\label{mht:eq:open-action}
 S_E(\mathcal A)=\int_U\tau_E\operatorname{tr}
 \left(\frac12\mathcal A(d_U+d_E)\mathcal A+\frac13\mathcal A^3\right).
\end{equation}
Only surface degree two and internal degree one contribute. On a
noncompact surface take compactly supported fields.

For the residual fields, write
$a\in\Omega^1(U)\otimes\operatorname{Mat}_N(\mathbb C)$ and
$J\in C^\infty(U)\otimes\operatorname{Mat}_N(\mathbb C)$.
Their image is $\mathcal A=a+J\eta_0$.

\begin{proposition}\label{mht:prop:bf-density}
With $F_a=d_Ua+a\wedge a$, the density of
\eqref{mht:eq:open-action} is
\[
 \kappa\operatorname{tr}(JF_a)
       -\frac{\kappa}{2}d_U\operatorname{tr}(aJ).
\]
On a closed surface this is the BF action
$\kappa\int_U\operatorname{tr}(JF_a)$. If $U$ has boundary, the
additional term is
$-\frac{\kappa}{2}\int_{\partial U}\operatorname{tr}(aJ)$.
\end{proposition}

\begin{proof}
The kinetic term has trace
\[
 \frac{\kappa}{2}\operatorname{tr}(a\wedge d_UJ+Jd_Ua)
 =\kappa\operatorname{tr}(Jd_Ua)
       -\frac{\kappa}{2}d_U\operatorname{tr}(aJ).
\]
The three terms of $\mathcal A^3$ containing one internal form are
$a^2J\eta_0$, $-aJa\eta_0$, and $Ja^2\eta_0$. Graded cyclicity makes
their traces equal, so the cubic term contributes
$\kappa\operatorname{tr}(Ja^2)$. All other terms vanish by degree.
Stokes' theorem gives the boundary integral.
\end{proof}

The inclusion now identifies states, products, trace, and the action on
the chosen residual fields. To evaluate an invariant on an arbitrary
smooth connection, we still need its residual matrix. The linear
projection need not compute it: the gauge transformation that eliminates
the other modes is nonlinear.
